\documentclass[11pt,a4paper]{article}
\usepackage{fontspec}\usepackage{polyglossia}
\setdefaultlanguage{russian}\setotherlanguage{english}
\newcommand{\optserif}{Path=fonts/,Extension=.ttf,UprightFont=*-Regular,BoldFont=*-Bold,ItalicFont=*-Italic,BoldItalicFont=*-BoldItalic,Ligatures=TeX}
\IfFileExists{fonts/LiberationSerif-Regular.ttf}{%
  \setmainfont{LiberationSerif}[\optserif]\newfontfamily\cyrillicfont{LiberationSerif}[\optserif]}{%
  \IfFontExistsTF{Liberation Serif}{\setmainfont{Liberation Serif}\newfontfamily\cyrillicfont{Liberation Serif}}{%
    \setmainfont{DejaVu Serif}\newfontfamily\cyrillicfont{DejaVu Serif}}}
\usepackage[a4paper,margin=1.9cm]{geometry}
\usepackage{enumitem,amsmath,booktabs,titlesec,array,longtable}
\usepackage[hidelinks]{hyperref}
\setlist[enumerate]{leftmargin=1.4em,itemsep=2pt,topsep=2pt}
\setlength{\parskip}{2pt}\emergencystretch=2em
\titleformat{\section}{\large\bfseries}{\thesection.}{0.4em}{}
\titlespacing*{\section}{0pt}{8pt plus 1pt}{3pt}
\begin{document}
\begin{center}
{\large\bfseries Что модель предсказывает измеримого}\\[2pt]
{\small Приложение к отчёту. 2 ноября 2026 г.}
\end{center}
\vspace{-2pt}

\noindent\textbf{Зачем этот список.} Вся работа сделана \emph{in silico}, и первый вопрос к ней ---
что отсюда можно измерить. Ниже собрано всё, что модель утверждает о величинах, доступных опыту:
число, направление, и чем именно утверждение опровергается. Предсказания упорядочены по силе
проверки, а не по важности механизма.

\noindent\textbf{Оговорка об единицах, обязательная.} Сравнимы \emph{направления, порядки и
отношения}, а не абсолютные значения. Параметры редокс-ядра взяты из кинетики клеточных линий, а не
PV-клеток; величины переключателя сна выражены в милливольтах его собственной калибровки. Там, где
ниже стоит абсолютное число, оно приведено как величина эффекта внутри модели, и переносить его на
опыт без пересчёта нельзя.

\section{Полоса спасённых: немонотонность повреждения по уязвимости}
\textbf{Утверждение.} Популяция живёт в трёх фазах, а не в двух. Между здоровым ядром и рухнувшим
лежит полоса клеток, которые \emph{уязвимее} своих рухнувших соседей и при этом целы: доля
$0{,}00732\pm0{,}00042$ (семнадцать стандартных ошибок от нуля), на мелкой сетке около $8{,}5\,\%$
популяции. При замороженной медленной защите полоса обращается в \emph{ровный нуль}, при прайминге
расширяется в $1{,}82$ раза.

\textbf{Что измерять.} Повреждение PV-клеток не должно быть монотонной функцией их исходной
уязвимости. В ткани обязана находиться подгруппа клеток с высокими показателями уязвимости и
сохранным парвальбумином, соседствующая с повреждёнными. Прокси уязвимости: целость перинейронной
сети, митохондриальные маркеры, показатели окислительной нагрузки на клетку.

\textbf{Чем опровергается.} Монотонной связью «уязвимее --- повреждённее» на одноклеточном уровне.
Это самое отличительное предсказание работы: оно контринтуитивно, не следует ни из какой усреднённой
картины и требует именно одноклеточного, а не тканевого разрешения.

\section{Предвестник в сне: ухудшение до срыва, и его источник}
\textbf{Утверждение.} Сон портится \emph{до} порога срыва, и делает это медленная память защиты, а
не медленный шум нагрузки. При отнятии медленной памяти высота предвестника падает у 107 наборов из
174, $p=0{,}0015$. Высота в последние сутки перед срывом: $1{,}061$ от собственного здорового уровня
набора, при аналитической вилке переключателя $1{,}0500$--$1{,}1280$.

\textbf{Что измерять.} Удлинение бодрствования порядка нескольких процентов в сутки,
предшествующие эпизоду. Данных такого рода много: актиграфия и полисомнография в когортах
клинического высокого риска и при шизофрении, где нарушения сна перед обострением описаны.

\textbf{Чем опровергается.} Отсутствием систематического сдвига перед эпизодом; либо наличием
сдвига, который не исчезает при вмешательстве, подавляющем медленное звено защиты. Второе --- более
сильная проверка, но требует животной модели.

\section{Потолок пассивной починки и аффинный закон длительности}
\textbf{Утверждение.} Длительность эпизода задаёт \emph{память} защиты, а не её мощность: внутри
набора $99{,}8\,\%$ дисперсии длительности объясняется ходом нагрузки и $0{,}2\,\%$ --- починкой
(разложение точное). Закон аффинный и не насыщается, длительность $=1{,}28+0{,}97\,\tau_b$ суток;
насыщающаяся форма с общим сдвигом отвергнута. Отсюда следствие: \emph{пассивное усиление починки
клинической цели не достигает} --- промах в 2--10 раз, поскольку нужная память около 12\,ч, а всякое
измеренное значение второй ступени лежит в 24--88\,ч.

\textbf{Что измерять.} Вмешательства, удлиняющие или укорачивающие время жизни защитного белка,
должны сдвигать длительность эпизода \emph{линейно}, без насыщения. Вмешательства, повышающие
мощность починки при неизменной её памяти, не должны сдвигать её вовсе.

\textbf{Чем опровергается.} Насыщением отклика при больших $\tau_b$; либо заметным сокращением
эпизодов от усиления починки без изменения её памяти.

\section{Прайминг работает верхом окна, а не скоростью починки}
\textbf{Утверждение.} Предварительная мягкая нагрузка укорачивает последующий эпизод: односторонняя
95-процентная граница медианы длительности $1{,}51$\,сут, внутри клинических двух суток. Механизм
назван: прайминг поднимает \emph{верх} окна, а не ускоряет починку. Количественно: ускорение первой
ступени индукции даёт от $-1$ до $-5\,\%$ длительности, усиление --- $-52\,\%$.

\textbf{Что измерять.} Гормезисный протокол: мягкое окислительное воздействие до основного должно
укорачивать эпизод. Разведение двух механизмов --- вмешательство, ускоряющее индукцию, против
вмешательства, усиливающего её потолок; модель предсказывает, что работает только второе.

\textbf{Чем опровергается.} Равной эффективностью обоих вмешательств либо преимуществом ускорения.

\section{Петля ёмкости: эпизоды реже, но длиннее}
\textbf{Утверждение.} Поклеточная потеря ёмкости защиты не учащает срывы, а удлиняет их:
длительность $13{,}1\to25{,}6$\,сут при падении частоты $0{,}042\to0{,}033$ на набор-сутки под
риском; эпизодичность сохраняется на уровне $0{,}75$ от разомкнутой модели. Прогрессирование
превышает сопоставленную по ёмкости статику на $+0{,}0231$ при MDE $0{,}0096$, воспроизведено на
трёх независимых реализациях шума.

\textbf{Что измерять.} В продольных когортах: по мере накопления повреждения эпизоды должны
становиться \emph{реже и длиннее}, а не чаще. Это проверяется на уже собираемых данных о числе и
длительности обострений.

\textbf{Чем опровергается.} Ростом частоты обострений при прогрессировании. Оговорка: абсолютное
значение $25{,}6$\,сут лежит много выше клинических двух суток, то есть рабочая точка петли
завышена; сравнимо здесь направление и отношение, не величина.

\section{Гистерезис: удержание зависит от глубины снятия нагрузки}
\textbf{Утверждение.} Сломанное состояние удерживается после снятия нагрузки, и удержание растёт
с мягкостью снятия монотонно ($\rho=+1{,}000$, перестановочный $p=0{,}0083$): при самом мягком
снятии $0{,}561$ против $0{,}146$ у разомкнутой модели.

\textbf{Что измерять.} Вероятность возврата после отмены лечения должна монотонно зависеть от
полноты снятия нагрузки, а не быть постоянной. Проверяемо в исследованиях отмены, хотя «полнота
снятия» требует отдельной операционализации.

\section{Рваная нагрузка не бьёт сильнее гладкой}
\textbf{Утверждение.} При закреплённых среднем и дисперсии частота срывов \emph{падает} на
$40\,\%$ с ростом времени корреляции нагрузки. Редкие сильные толчки не опаснее ровного подъёма ---
гипотеза, которую мы проверяли и опровергли.

\textbf{Что измерять.} Протоколы стресса с одинаковым средним и дисперсией, но разной временной
структурой: сгруппированный стресс должен давать \emph{меньше} эпизодов, чем равномерный.
Проверяемо на животной модели.

\section{Серотониновый всплеск --- ручка громкости, а не спусковой крючок}
\textbf{Утверждение.} За всплеском следует $26\,\%$ срывов при аналитическом ожидании $43\,\%$ и
предрегистрированном пороге $80\,\%$; ни одно свойство эпизода при этом не меняется.

\textbf{Что измерять.} Серотонинергическая провокация не должна надёжно запускать эпизод и не должна
менять его длительность или глубину.

\section{Асимметрия защиты задаёт её средний уровень}
\textbf{Утверждение.} Если постоянные подъёма и спада защитного белка различаются в $r$ раз, средний
уровень защиты стоит в точке $w(r)=fr/(fr+1-f)$ при $f=0{,}705$ (измеренная доля бодрствования).
Перебор $r=0{,}33$--$3$ предсказывает $w=0{,}441$--$0{,}878$; измеренные уровни защиты следуют этому
порядку точно ($\rho=+1$, $p=0{,}0083$). В формуле нет ни одного подогнанного числа.

\textbf{Что измерять.} Обе постоянные защитного белка порознь, в клетке. Модель предсказывает
средний уровень защиты без свободных параметров.

\section{Два отрицательных предсказания}
\textbf{Нагрузочная петля не может дать хроничность и эпизодичность одновременно.} Необратимость
включается при силе петли $1{,}11$, эпизоды гаснут при $0{,}18$ --- шестикратный структурный разрыв,
устойчивый к более мягкому протоколу. Следствие для опыта: если у больных надёжно наблюдаются
\emph{и} необратимое накопление, \emph{и} сохранная эпизодичность, то механизм не в петле через
общую нагрузку. Причина названа: нагрузка одна на набор, поэтому петля через неё двигает популяцию
целиком.

\textbf{Поломка ритма лестницей захвата недостижима при физиологической связи.} Первая ступень
требует $k\!\cdot\!D=4{,}77$\,мВ, а тождественно $k\!\cdot\!D\le k$, и при рабочем $k=2$\,мВ она не
достигается ни при каком повреждении. Модель, таким образом, предсказывает, что циклов сна в 48 и
72\,ч при шизофрении быть не должно --- и это согласуется с тем, что их и не описывают. Предсказание
слабое, но оно единственное в работе, где несовпадение модели с опытом оказалось в её пользу.

\section{Чего модель не предсказывает}
Величины внутреннего устройства --- отношение границ окна, положения складок в милливольтах,
положение нагрузки в окне, число клеток в наборе --- не являются предсказаниями: они либо выбраны,
либо служебны. Дофамин и позитивные симптомы вне модели по замыслу. И главное: модель описывает
возникновение и прекращение \emph{эпизода}, а не течение болезни, поэтому любое утверждение о
масштабе лет лежит вне её области применимости.
\end{document}
